\vspace{+1em}

\begin{center}
\large\textcolor{Red}{\textbf{Preface}}
\end{center}

% \begin{multicols*}{2}

\ubsection{Preface}{preface-breviarium}

% {\color{Red} \ubsubsection{I. Historical Background}{historical-background}}

\lettrine[lines=2]{\zallmancaps \color{Blue} T}{he} fame of the Prototype of Humbert precedes its contents, so a brief introduction suffices. As the story goes, and as highlighted by Fr. William R. Bonniwell, O.P. (\emph{A History of the Dominican Breviary}, New York 1945), in the 13th century, the newly-organized Dominicans, as an international order, wished to possess a uniform liturgy across their houses. After various attempts, the fifth Master-General, Humbert of Romans, commissioned a new standard and imposed it on the order, which was so confirmed by Pope Clement IV in 1267, at the request of Bl. John of Vercelli, the succeeding Master-General.

The Codex, \emph{Santa Sabina XIV L1}, is held at Sants Sabina in Rome, and is divided across fourteen books:
% \begin{itemize}[noitemsep]
\begin{itemize}[nosep]
	\item \emph{Ordinarium} (Ordinal)
	\item \emph{Martyrologium} (Martyrology)
	\item \emph{Collectarium} (Collectary)
	\item \emph{Processionarium} (Processionary)
	\item \emph{Psalterium} (Psalter)
	\item \emph{Breviarium} (Breviary)
	\item \emph{Lectionarium} (Lectionary)
	\item \emph{Antiphonarium} (Antiphonary)
	\item \emph{Graduale} (Gradual)
	\item \emph{Pulpitarium} (Pulpitary)
	\item \emph{Missale Conventuale} (Missal for Sung Mass)
	\item \emph{Epistolarium} (Epistolary)
	\item \emph{Evangelistarium} (Evangelistary)
	\item \emph{Missale Minorum Altarium} (Missal for Low Mass)
\end{itemize}

Of these, only a precious few books have been transcribed:
\begin{itemize}[nosep]
	\item The \emph{Ordinarium} (and small selections of other books) was transcribed in 1921 by Fr. Franciscus M. Guerrini, O.P. (\emph{Ordinarium Juxta Ritum Sacri Ordinis Fratrum Praedicatorum}, Rome 1921).
	\item The Sanctoral half of the \emph{Lectionarium} in 1995 by the late Anne-Élisabeth Urfels-Capot, PhD. (\emph{Le sanctoral du lectionnaire de l'office dominicain (1254-1256)}, Paris 2007).
	\item Regarding the \emph{Antiphonarium}, the late Holger Peter Sandhofe transcribed a portion of these chants before his untimely death in 2005.
	\item Recently, in 2021, Fr. Dominik Jurczak, O.P., SLD., transcribed the \emph{Evangelistarium} and the Gospel incipits from the \emph{Martyrologium} (\emph{Vangelo e la liturgia domenicana dopo la riforma di Umberto di Romans}, Rome 2021).
\end{itemize}

This Breviary presents the entirety of the \emph{Breviarium}, the text of the \emph{Psalterium} (minus the chant notation and Psalm tones), and the relevant parts of the \emph{Martyrologium} and \emph{Collectarium} to facilitate recitation of the Divine Office according to Humbert.


% https://institutumhistoricum.op.org/vol-38-vangelo-e-la-liturgia-domenicana-dopo-la-riforma-di-umberto-di-romans/
% https://institutumhistoricum.op.org/en/vol-38-vangelo-e-la-liturgia-domenicana-dopo-la-riforma-di-umberto-di-romans/
% https://angelicum.it/publication/vangelo-e-la-liturgia-domenicana-dopo-la-riforma-di-umberto-di-romans-con-ledizione-dellevangelistarium-e-dellincipit-evangeliorum-dal-codice-xiv-l1-di-santa-sabina/
% https://libreriasoleeditrice.com/products/vangelo-e-liturgia-domenicana-dopo-la-riforma-di-umberto-di-romans-con-l-edizione-dell-incipit-evangeliorum-e-dell-evangelistarium-dal-codice-xiv-l1-di-santa-sabina-9788899616380
	% €90,00
% https://www.ibs.it/vangelo-liturgia-domenicana-dopo-riforma-libro-dominik-jurczak/e/9788899616380
	% 85,50 €
% https://www.unilibro.it/libro/jurczak-dominik/vangelo-liturgia-domenicana-dopo-riforma-umberto-romans-edizione-dell-incipit-evangeliorum-evangelistarium-codice-xiv-l1-santa-sabina/9788899616380
	% € 85.50
% https://studium.op.org/vangelo-e-liturgia-domenicana-dopo-la-riforma-di-umberto-di-romans/
% https://shop.librerialenuvole.it/scheda-libro/dominik-jurczak/vangelo-e-liturgia-domenicana-dopo-la-riforma-di-umberto-di-romans-con-ledizione-dellincipit-evangeliorum-e-dellevangelistarium-dal-codice-xiv-l1-di-santa-sabina-9788899616380-3775526.html
	% 90,00 €
% https://www.libreriauniversitaria.it/vangelo-liturgia-domenicana-dopo-riforma/libro/9788899616380
	% € 90,00
% https://www.libreriadelsanto.it/libri/9788899616380/vangelo-e-liturgia-domenicana-dopo-la-riforma-di-umberto-di-romans.html
	% 85,50 €




% elec.enc.sorbonne.fr/sanctoral/sanctoral


% https://www.amazon.com/-/he/Anne-Elisabeth-Urfels-Capot/dp/2900791928
% https://www.amazon.fr/Aux-origines-liturgie-dominicaine-Manuscrit/dp/2271056322
% https://www.amazon.fr/Aux-origines-liturgie-dominicaine-manuscrit/dp/2728305528
% https://www.torrossa.com/en/resources/an/2250867
% https://www.cnrseditions.fr/catalogue/histoire/aux-origines-de-la-liturgie-dominicaine-le-manuscrit-santa-sabina-xiv-l-1/
% https://www.persee.fr/issue/dirht_0073-8212_2004_act_67_1
% https://books.google.com/books/about/Aux_origines_de_la_liturgie_dominicaine.html?id=y-URAvZccYkC



{\color{Red} \ubsubsection{Rationale \& \ Problem Statement}{rationale-problem-statement}}

\par Since the dawn of Christianity, Catholics have prayed the canonical hours and Psalms for two millennia, continuing the ancient Jewish Temple tradition, and our current time is no exception. From the height of the Middle Ages, when every order and secular cathedral sang its own Use of the Office, to laymen reciting various \textit{Libri Horarum}, and the rise of the printing press and the Tridentine standard, the number of local Uses has steadily decreased in popularity, leaving the Dominicans, along with a handful of other sees and orders, retaining their own Use.

Catholics today wishing to pray the hours have a plethora of options, yet a strange dearth of the traditional variety. One can purchase various editions of the \textit{nouveau} Liturgy of the Hours, as well as various reprints of the 1962 secular and 1963 monastic breviaries, yet outside of those, we venture into the vernacular with editions of the same, with the added choice of various Eastern Offices, from Byzantine, to Eastern Syriac, to Coptic, to Armenian. Yet within the Roman Rite itself, we are left in a strange situation of being content with the new, reformed books of Vatican II, reprints of the 1962 books, or we are left outbidding each other on online auctions for ancient Breviaries last printed seventy or eighty years ago. The same holds even more true and dire for those wishing to bypass the Breviary reforms of Pius X.

{\color{Red} \ubsubsection{Vision \& Editorial Principles}{vision-and-editorial-principles}}

% II. The Project: Vision \& Editorial Principles

What I present here is a labor of two years in the making, starting in June 2024 with these several goals:
\begin{itemize}[nosep]
% \begin{itemize}[nosep,leftmargin=8pt]
	\item To typeset a Pre-\textit{Divino Afflatu} Breviary,
	\item To recover the Dominican Rite, as standardized by Bl. Humbert of Romans,
	\item To show \LaTeX \ is viable for typesetting Breviaries, providing an open-source template, and
	\item To demonstrate this using modern conveniences, such as Print-on-Demand services.
\end{itemize}

Thus, by the grace of God, this is the first \emph{Totum}, Dominican Rite, Latin Pre-55, and Pre-\textit{Divino Afflatu} Breviary printed in more than 60 years.

I have endeavoured to keep the layout of this Breviary as faithfully as possible to the 13th-century manuscript of Humbert, namely, the Prototype of Santa Sabina XIV L1, which was the foundational text that unified early Dominicans in prayer. This is the Rite that Sts. Thomas Aquinas, Albert the Great, and Raymond of Penyafort prayed, which has been restored and newly typeset, without the cluttered Kalendar that is characteristic of later editions.

By bypassing modern and Baroque accretions, this text allows contemporary laity, religious, and scholars to experience the Divine Office as it was laid down by Bl. Humbert, and as it was prayed by early Dominican saints. This Breviary bridges the gap between historical fidelity and accessibility for daily prayer.

For the Order of Preachers, the solemn celebration of the Divine Office is their primary duty and the source of their contemplative and doctrinal life. May this edition serve to nourish your own contemplation, that the Church may bear her fruits to the world.


% \end{multicols*}